\begin{prob}
    Determine the time complexity of the following piece of code, showing your reasoning
    and your work.

    \inputminted{python}{\thisdir/include/code.py}

    \textbf{Hint}: you might need to think back to calculus to remember the formula for
    the sum of a geometric progression... or you can check wikipedia.%
    \footnote{\url{https://en.wikipedia.org/wiki/Geometric_progression}}

    \begin{soln}
        $\Theta(n)$.

        How many times does the \python{print} statement execute? That will tell us the time complexity.
        On the first iteration
        of the outer loop, \python{i = 2} and this line runs twice. On the second iteration
        of the outer loop, \python{i = 4} and this line runs four times. On the third
        iteration, it runs 8 times. In general, on the $k$th iteration of the outer loop,
        the line runs $2^k$ times.

        The outer loop will run \emph{roughly} $\log_2 n$ times, since $i$ is doubling on
        each iteration. We'll see in a moment why it's OK to have a rough
        estimate instead of being exact here.

        The total number of executions of the \python{print} is therefore:
        \[
            \underbrace{2}_{\text{\footnotesize 1st outer iter.}}
            +
            \underbrace{4}_{\text{\footnotesize 2nd outer iter.}}
            +
            \underbrace{8}_{\text{\footnotesize 3rd outer iter.}}
            +
            \ldots
            +
            \underbrace{2^k}_{\text{\footnotesize $k$th outer iter.}}
            +
            \ldots
            +
            \underbrace{2^{\log_2 n}}_{\text{\footnotesize $\log_2 n$th outer iter.}}
        \]

        Or, written using summation notation:
        \[
            \sum_{k=1}^{\log_2 n} 2^k
        \]

        This is the sum of a \emph{geometric progression}. Wikipedia tells us that the formula
        for the sum of $1 + 2 + 4 + 8 + \ldots + 2^K$ is:
        \[
            \sum_{k=0}^{K} 2^k = \frac{1 - 2^{K+1}}{1 - 2} = 2^{K+1} - 1
        \]
        Notice that this sum starts from $k = 0$, while ours starts with $k = 1$. 
        In other words, our sum is the same except it is missing the first term corresponding
        to $k = 0$.
        So to 
        ``correct'' this, we subtract the $k = 0$ term (which is $2^0 = 1$)
        from the larger sum :
        \[
            \sum_{k=1}^{K} 2^k = 
            \left(\sum_{k=0}^{K} 2^k\right) - 2^0 = 
            \left(2^{K+1} - 1\right) - 1=
            2^{K+1} - 2
        \]
        Plugging in $K = \log_2 n$:
        \begin{align*}
            \sum_{k=1}^{\log_2 n} 2^k 
            &= 2^{\log_2 n +1} - 2
            \intertext{Using the fact that $a^{b+c} = a^b \cdot a^c$:}
            &=
                2^{\log_2 n} \cdot 2 - 2\\
            &=
                2n - 2\\
            &=
                \Theta(n)
        \end{align*}

        Now, remember that we said that we could afford to be a little imprecise when calculating
        the number of times that the outer loop runs. Let's see why. If $n$ isn't a power
        of $2$, $\log_2 n$ will not be an integer. For instance, if $n = 17$, $\log_2 n = 4.08$.
        Of course, the loop can't iterate 4.08 times -- you can check that it will actually iterate
        5 times. In general, the loop will run exactly $\lceil \log_2 n \rceil$ times, where
        $\lceil \cdot \rceil$ is the \emph{ceiling} operation; it rounds a real number up
        to the next integer.

        In other words, $\log_2 n$ could actually be as much as one less than the actual number
        of iterations. Perhaps to be careful we should overestimate the number of iterations
        to be $\log_2 n + 1$. What if we were to use $\log_2 n + 1$ instead? We'd end up with:
        \[
            \sum_{k=1}^{\log_2 n + 1} 2^k = 2^{\log_2 n + 2} - 2 = 4\cdot 2^{\log_2{n}} - 2 = \Theta(n) = 4n-2
        \]
        So the time complexity doesn't change. This is why using $\Theta$ notation allows us
        to be ``sloppy'' at times without being incorrect. It saves us work, as long as
        we know how to use it correctly!
    \end{soln}

\end{prob}
